\begin{afterword}
\summaryhead

The post starts from typicality. Most people judge a robin a more typical bird than an ostrich, and a desk chair a more typical chair than a beanbag. Psychologists measure this in several ways that agree, such as how fast people confirm that a robin or a penguin is a bird, and direct ratings.

The author then looks for a bias that typicality produces. Judgments of similarity are asymmetric: ``98 is approximately 100'' sounds more natural than the reverse, and people rate Mexico as more similar to the United States than the United States to Mexico. In a study by Rips, people expected a disease to spread from robins to ducks more readily than from ducks to robins. The author argues that no practical difference could justify this, and explains it by an asymmetric sense of similarity used in place of probability. A spatial example follows: people seem to treat Kansas as close to everything and Alaska as far from everything. The post concludes that the Aristotelian model of categories fails, and that people treat membership in a category as a matter of degree.

\responsehead

The post reports the direction of its findings correctly. Typical members are verified faster, as the cited volume says. The country ratings in Tversky and Gati's table go the way the post says, 7.65 against 6.46 for Mexico and the United States. Across all their pairs the effect was small, 0.42 points on a scale of 1 to 20, and six of the 21 pairs went the other way.

The sources for the first example need correcting. The comparison of 98 and 100 is cited to a paper by Sadock that does not contain it. The experiment that supports it is Rosch's work on reference points, reported in the volume the post cites a paragraph earlier. The point survives the correction.

The central example is the disease study, and there the post goes further than the evidence on two counts. First, it says that in a practical sense any difference that would slow a disease from ducks to robins must slow it just as much from robins to ducks. A published Bayesian model (Heit 1998) disagrees. If atypical species are believed to have many features of their own, a new property found in ducks is more likely to belong to ducks alone. Second, the post's explanation, an asymmetric similarity heuristic, is not the one the study supports. Rips measured similarity symmetrically and found a separate effect of the typicality of the species in the premise; the asymmetry follows from that. Both accounts make typicality central, so the post's topic survives. Its mechanism does not follow from the study, and its verdict that the answers have no practical basis is disputed.

The conclusion also goes beyond the evidence. That people rate some members of a category as more typical does not show that they treat membership as a matter of degree. People rate 3 as a more typical odd number than 4,284,647, yet both are ``definitely odd, not vaguely so,'' in the words of Rips and a coauthor. The Aristotelian model may fail for other reasons. Typicality alone does not show that ``Statements of set membership can be more or less true.''

In short: accurate reports of typicality and asymmetry; a disease study whose asymmetry the study itself and later models explain by typicality, not by the asymmetric similarity the post proposes; and a conclusion about graded membership that typicality does not support.
\end{afterword}
